\section{Materiais e Métodos}

\subsection{Visão geral do pipeline}

A Figura~\ref{fig:pipeline-v2} resume o pipeline em seis etapas: (E1) lista de municípios; (E2) descoberta de instâncias do SAPL; (E3) extração de PLs; (E4) tradução \textit{ementa} $\rightarrow$ \textit{ação}; (E5) indexação vetorial; e (E6) análise com indicadores e agrupamento.

\begin{figure*}[t]
\centering
\resizebox{\textwidth}{!}{%
\begin{tikzpicture}[node distance=0.9cm, font=\small]
\tikzstyle{box}=[draw, rounded corners, align=center, inner sep=6pt]
\node[box] (e1) {E1\\Lista de\\municípios\\(IBGE)};
\node[box, right=0.6cm of e1] (e2) {E2\\Descoberta\\de instâncias\\do SAPL};
\node[box, right=0.6cm of e2] (e3) {E3\\Extração\\de PLs};
\node[box, right=0.6cm of e3] (e4) {E4\\Tradução\\de ementa\\em ação};
\node[box, right=0.6cm of e4] (e5) {E5\\Indexação\\vetorial\\(E5 + Qdrant)};
\node[box, right=0.6cm of e5] (e6) {E6\\Indicadores\\oficiais\\+ Agrupamento};
\draw[->] (e1) -- (e2);
\draw[->] (e2) -- (e3);
\draw[->] (e3) -- (e4);
\draw[->] (e4) -- (e5);
\draw[->] (e5) -- (e6);
\end{tikzpicture}
}%
\caption{Pipeline conceitual do Sonar Municipal na versão reestruturada do artigo.}
\label{fig:pipeline-v2}
\end{figure*}

\subsection{Descoberta de instâncias do SAPL}

O universo inicial foi definido pela lista de municípios da API de localidades do IBGE \cite{ibge_localidades_api}, reduzindo viés de seleção por listas manuais. A descoberta combina heurísticas baseadas em nomes de municípios e validação por evidência pública. O \textit{finder} gera candidatos como \texttt{sapl.\{slug\}.\{uf\}.leg.br} e \texttt{\{slug\}.\{uf\}.leg.br}, testa rotas típicas do SAPL e confirma a instância quando a resposta contém marcadores compatíveis ou título de página correspondente.

Neste trabalho, consideramos uma \textit{instância SAPL validada} como aquela que passou pelas heurísticas de descoberta e pela checagem pública de compatibilidade. A implementação de referência deduplica resultados por URL validada e grava o progresso em JSONL, o que favorece auditoria e retomada de execução.

\begin{algorithm}[H]
\caption{Descoberta de instâncias SAPL}
\label{alg:finder-v2}
\begin{algorithmic}[1]
\Require Lista de municípios $M$
\Ensure Instâncias SAPL validadas $S$
\State $S \gets \emptyset$
\For{cada município $m \in M$}
  \State gerar candidatos de host a partir do nome e da UF de $m$
  \For{cada candidato $c$}
    \If{$c$ expõe rota pública com marcadores compatíveis com SAPL}
      \State validar $c$ e adicionar em $S$
      \State interromper busca daquele município
    \EndIf
  \EndFor
\EndFor
\State deduplicar $S$ por URL validada
\State \Return $S$
\end{algorithmic}
\end{algorithm}

\subsection{Extração de Projetos de Lei}

Uma vez validada a instância, a extração deriva a base correta do SAPL, testa os endpoints de matéria legislativa e recupera o catálogo de tipos de matéria. Todos os tipos cujo rótulo normalizado contém ``projeto de lei'' são tratados como candidatos válidos. A extração percorre a paginação até esgotar os resultados e deduplica itens por origem da instância e identificador da matéria.

Neste artigo, definimos \textit{extração bem-sucedida} como a obtenção de pelo menos um PL com metadados mínimos válidos. A implementação também aplica deduplicação por município, número do PL, ano e origem da instância.

\begin{algorithm}[H]
\caption{Extração exaustiva de PLs}
\label{alg:scrapper-v2}
\begin{algorithmic}[1]
\Require Instância SAPL validada $s$
\Ensure Conjunto de PLs $P$
\State derivar a base da API de $s$
\State consultar tipos de matéria disponíveis
\State selecionar tipos com rótulo compatível com ``Projeto de Lei''
\For{cada tipo selecionado}
  \State percorrer todas as páginas do endpoint de matérias
  \State recuperar metadados relevantes para cada PL
\EndFor
\State deduplicar $P$ por origem e identificadores legislativos
\State \Return $P$
\end{algorithmic}
\end{algorithm}

\subsection{Snapshot experimental e cobertura}

Os números reportados neste artigo correspondem a uma execução congelada em 12/11/2025. Nessa execução, foram encontradas 1.259 instâncias do SAPL e 322 instâncias permitiram extração automatizada com ao menos um PL válido, totalizando 220.065 PLs.

Tomando os 5.570 municípios brasileiros como referência, a cobertura foi de aproximadamente 22,60\% para descoberta de instâncias e 5,78\% para extração automatizada com sucesso. O dataset final cobre 322 municípios em 19 UFs. Essa discrepância entre descoberta e extração reflete heterogeneidade técnica entre portais, mudanças locais de rota, proxies, autenticação, variações nos endpoints e barreiras operacionais ao acesso programático. Foi realizada validação manual exploratória, sem desenho estatístico formal, apenas como \textit{sanity check}.

\subsection{Qualidade estrutural do dataset final}

O dataset final contém 220.065 registros, cada um com metadados legislativos, ementa, ação textualizada e embedding de dimensão 768. Todos os registros possuem \textit{ementa}, \textit{ação}, \textit{link\_publico} e embedding preenchidos. O campo \textit{situacao}, por outro lado, está ausente em toda a base e, portanto, não foi usado nas análises desta versão.

Do ponto de vista de integridade, a base final não apresentou duplicatas por \textit{link\_publico} nem por par \texttt{(sapl\_base, materia\_id)}, o que sugere deduplicação consistente no nível do item legislativo público. Ao mesmo tempo, a inspeção exploratória mostrou que o critério de filtragem por rótulo textual não produz uma base composta exclusivamente por PLs estritos: cerca de 3,54\% dos registros correspondem a anteprojetos, emendas, substitutivos ou vetos cujos rótulos ainda contêm a expressão ``projeto de lei''. Essa contaminação é pequena em termos relativos, mas metodologicamente relevante.

\subsection{Tradução de ementas em ações}

As ementas dos PLs são curtas, variáveis e fortemente marcadas por linguagem jurídica. Para melhorar a interpretabilidade e reduzir ruído na indexação, o projeto traduz cada ementa em uma ação operacional, no formato ``verbo no infinitivo + objeto + complementos essenciais''. A ideia é representar o núcleo material da política, e não sua casca normativa.

O conjunto de treinamento foi composto por 1.000 ementas reais distintas, selecionadas aleatoriamente. Para cada ementa, gerou-se uma ação correspondente por meio de um prompt instrucional aplicado ao GPT-5.1-Thinking \cite{openai_gpt51_model,openai_gpt51_system_card}. O prompt utilizado está disponível no repositório do projeto. Após a geração, foram corrigidos casos com perda de sentido, excesso de abstração e resíduos de juridiquês.

Como modelo base, adotou-se o PTT5-v2 \cite{piau_ptt5v2_2024}, uma variante do T5 \cite{raffel_t5_2020} treinada para português. O ajuste fino foi realizado com QLoRA em 4 bits \cite{dettmers_qlora_2023}. O conjunto foi dividido em 90/10 para treino e validação. O treinamento foi executado por 30 épocas, com taxa de aprendizado de $3\times 10^{-4}$ e \textit{batch size} 16. A avaliação utilizou BERTScore \cite{zhang_bertscore_2019}, obtendo 84\% no conjunto de validação.

No acervo completo, a tradução reduz a extensão média do texto de 152,86 para 90,01 caracteres e de 24,18 para 13,61 palavras, o que corresponde a aproximadamente 58,9\% dos caracteres e 56,3\% das palavras da ementa original. Além de encurtar o texto, a tradução aumenta a convergência textual: enquanto as ementas têm apenas 0,37\% de reuso exato, as ações apresentam 14,46\% de reuso, o que favorece busca e agrupamento.

O principal \textit{baseline} considerado até aqui foi não traduzir a ementa e indexar o texto legislativo bruto. Também houve revisão informal por especialistas, mas ainda sem protocolo quantitativo formal.

\subsection{Busca semântica}

As ações geradas são indexadas por embeddings com o modelo Multilingual E5 \cite{wang_e5_2024}. A implementação segue o padrão recomendado pelo modelo, usando prefixo \texttt{query:} nas consultas. As embeddings são armazenadas no Qdrant \cite{qdrant_overview,qdrant_collections}, e o sistema retorna os itens com maior similaridade vetorial.

A avaliação atual da recuperação é preliminar: utilizamos consultas reais e julgamento manual de relevância das ações retornadas nas posições mais altas. Ainda não foi construído um benchmark anotado formalmente, nem métricas como Precision@K, Recall@K ou nDCG foram consolidadas para publicação. Em termos operacionais, o valor de $K = 1000$ foi adotado empiricamente após testes exploratórios, priorizando diversidade de candidatos para a etapa posterior de agrupamento.

\subsection{Indicadores municipais e janelas de efeito}

Para aproximar o sistema de desfechos observáveis, utilizamos dois indicadores municipais oficiais: taxa de homicídios por 100 mil habitantes e taxa de matrículas em ensino regular por 100 mil habitantes. As fórmulas são:

\[
\text{Taxa de homicídios} = \frac{\text{homicídios}}{\text{população}} \times 100.000
\]

\[
\text{Taxa de matrículas} = \frac{\text{matrículas no ensino regular}}{\text{população}} \times 100.000
\]

No caso da educação, tratamos matrículas como uma aproximação operacional de permanência escolar. Essa proxy é imperfeita: mudanças de matrícula podem refletir demografia, migração ou migração entre redes pública e privada, e não evasão em sentido estrito.

Os efeitos são calculados comparando o valor do indicador na data mais próxima à apresentação do PL com o valor observado em janelas posteriores de 6, 12, 18, 24, 30 e 36 meses. Essa granularidade foi escolhida por compatibilidade prática com a disponibilidade de séries temporais municipais. Municípios com dados insuficientes são excluídos da análise daquele indicador. A análise é puramente descritiva e associativa, sem controle causal explícito para tendência prévia, sazonalidade ou choques externos.

\subsection{Agrupamento em políticas públicas}

PLs semanticamente próximos são agrupados em políticas públicas para reduzir sparsidade analítica. Na implementação atual, o agrupamento usa apenas o texto da ação derivada da ementa. A similaridade é calculada por Jaccard entre conjuntos de tokens normalizados, e o limiar foi definido empiricamente por inspeção qualitativa dos grupos. A implementação de referência utiliza limiar 0,75.

Sejam $n$ o número de municípios representados no grupo e $n_{\text{positivos}}$ o número de municípios com efeito considerado favorável ao objetivo do indicador. O escore de qualidade é:

\[
Q = \frac{n_{\text{positivos}}}{n} \times \frac{n}{n+1}
\]

Esse escore combina taxa de sucesso observada com penalização simples para grupos muito pequenos. Também realizamos inspeção qualitativa para verificar se os PLs agrupados representavam, de fato, uma mesma política pública.

\subsection{Arquitetura do sistema}

A plataforma Sonar Municipal foi desenvolvida em React com Next.js e hospedada na Vercel \cite{vercel_nextjs}. A arquitetura integra interface web, serviço de embeddings, banco vetorial e dados de indicadores. A Figura~\ref{fig:arquitetura-v2} representa essa visão técnica.

\begin{figure}[H]
\centering
\resizebox{0.95\columnwidth}{!}{%
\begin{tikzpicture}[node distance=0.55cm, font=\scriptsize]
\tikzstyle{box}=[draw, rounded corners, align=center, inner sep=5pt]
\node[box] (u) {Usuário};
\node[box, below=of u] (w) {Frontend\\Next.js};
\node[box, below left=0.7cm and 0.9cm of w] (e) {Serviço de\\embeddings};
\node[box, below=of w] (q) {Qdrant};
\node[box, below right=0.7cm and 0.9cm of w] (i) {Indicadores\\oficiais};
\draw[->] (u) -- (w);
\draw[->] (w) -- (e);
\draw[->] (w) -- (q);
\draw[->] (w) -- (i);
\draw[->] (e) -- (q);
\end{tikzpicture}
}%
\caption{Arquitetura técnica simplificada da plataforma.}
\label{fig:arquitetura-v2}
\end{figure}

\subsection{Reprodutibilidade}

O código do projeto, incluindo coleta, treinamento, indexação e interface, está disponível em \url{https://github.com/thiagoambiel/SonarMunicipal}. A Tabela~\ref{tab:reprodutibilidade} resume os principais artefatos experimentais desta versão.

\begin{table}[H]
\centering
\caption{Resumo de reprodutibilidade da versão 2 do artigo.}
\label{tab:reprodutibilidade}
\small
\begin{tabular}{p{0.42\linewidth} p{0.5\linewidth}}
\toprule
Item & Configuração reportada \\
\midrule
Snapshot de coleta & 12/11/2025 \\
Instâncias SAPL encontradas & 1.259 \\
Instâncias com extração válida & 322 \\
PLs coletados & 220.065 \\
Pares ementa $\rightarrow$ ação & 1.000 \\
Divisão do tradutor & 90/10 (treino/validação) \\
Modelo tradutor & PTT5-v2 + QLoRA \\
Treinamento do tradutor & 30 épocas, LR $3\times10^{-4}$, batch 16 \\
Modelo de recuperação & Multilingual E5 \\
Valor de $K$ na busca & 1000 \\
Janelas de efeito & 6, 12, 18, 24, 30, 36 meses \\
Agrupamento & Jaccard sobre ações textualizadas \\
\bottomrule
\end{tabular}
\end{table}
